\chapter{Interprocedural Analysis}
\label{ch:ch22-interprocedural-analysis}

\section{One Function Is Not the Program}
\label{sec:ch22-one-function}

Every analysis so far has worked inside a single procedure.
It assumes unknown values at entry, treats calls as barriers that destroy facts, and hopes the damage stays small.
Real programs defeat that hope constantly.
Constants flow through arguments, objects escape through parameters, and the hottest facts live one call away from where they are used.
A load that looks redundant becomes provably redundant only when the callee is known not to store to the address.
A branch that looks unpredictable folds once the caller passes a constant.
\term{Interprocedural analysis} reasons across procedure boundaries so that optimization sees the whole program instead of one function at a time.
The price is scope: the analyzer must handle many procedures, many call sites, recursion, separate compilation, and code it cannot see.
The reward is that summaries computed once per procedure answer questions at every call site cheaply.
This chapter builds the call graph, defines what flows across it, computes procedure summaries with a worklist, and explains the precision ladder from context insensitivity to full sensitivity.
It generalizes the dataflow of earlier chapters and supplies the summaries that Chapters~\ref{ch:ch20-alias-analysis} and~\ref{ch:ch21-memory-optimization} need at call boundaries.

\section{Call Graphs}
\label{sec:ch22-call-graphs}

A \term{call graph} has one node per procedure and one edge per possible caller-to-callee transfer.
Direct calls contribute exact edges: the callee is named in the IR, so the edge is certain.
Indirect calls through function pointers, virtual methods, or callbacks contribute possible edges for every procedure whose address could flow to the call site.
Resolving those edges needs pointer analysis from Chapter~\ref{ch:ch20-alias-analysis}, which creates a productive cycle: better points-to sets give a smaller call graph, and a smaller call graph gives cheaper interprocedural facts.
Listing~\ref{lst:ch22-calls} shows a program whose graph is small enough to draw completely.

\begin{irlisting}[caption={Three procedures and a main program with direct calls and one global store each.},label={lst:ch22-calls}]
h(): store g1 = 1
     ret
g(): call h()
     store g2 = 2
     ret
f(): call g()
     ret
main: call f()
      x = load g1
      y = load g2
      ret x + y
\end{irlisting}

The graph has the edges $main \to f$, $f \to g$, and $g \to h$, with no cycles and no indirect calls.
Recursion would add a cycle, mutual recursion a strongly connected component, and separate compilation would add unknown nodes for code outside the analyzed units.
Unknown code forces sound approximations: a call to an unseen procedure is treated as reading and writing everything it could reach.
Engineers call this the \enquote{worst case} summary, and it is the default every sound whole-program analysis needs.
That treatment is always safe and often imprecise, which is why whole-program builds and link-time analysis exist.
Figure~\ref{fig:ch22-call-graph} draws the graph for Listing~\ref{lst:ch22-calls}.
The analysis propagates facts along these edges in reverse topological order when the graph is acyclic, and iterates to a fixed point when cycles appear.

\begin{figure}
  \centering
  \begin{tikzpicture}[node distance=16mm and 22mm]
    \node[basic block] (ch22main) {main};
    \node[basic block,below=of ch22main] (ch22f) {f};
    \node[basic block,below=of ch22f] (ch22g) {g};
    \node[basic block,below=of ch22g] (ch22h) {h};
    \draw[cfg edge] (ch22main) -- (ch22f);
    \draw[cfg edge] (ch22f) -- (ch22g);
    \draw[cfg edge] (ch22g) -- (ch22h);
  \end{tikzpicture}
  \caption{Direct calls give an exact acyclic call graph, so summaries propagate bottom-up from the leaves.}
  \label{fig:ch22-call-graph}
\end{figure}

\section{Summaries: May-Modify and Friends}
\label{sec:ch22-summaries}

A \term{procedure summary} describes a procedure's observable behavior once, so every caller can reuse the description.
The most widely used summary is the \term{may-modify} set, also called may-mod: the globals and address-taken locations the procedure or anything it calls might write.
Its dual is the \term{may-reference} set of locations it might read.
Constant-propagation summaries map input facts to output facts, alias summaries describe how pointer arguments connect to results, and purity summaries record that a procedure touches nothing observable at all.
Summaries compose along call edges: the summary of a caller includes its own direct effects plus the summaries of its callees.
That composition rule turns interprocedural analysis into a fixed-point problem over the call graph.
Algorithm~\ref{alg:ch22-maymod} computes may-modify sets for globals with a worklist over callers.

\begin{algorithm}
  \caption{May-modify summaries over the call graph}\label{alg:ch22-maymod}
  \begin{pseudocode}
    \Require fixed sound call graph, conservative direct store sets $Direct(P)$ including external effects
    \Ensure least fixed point of may-modify sets $Mod(P)$
    \State $Mod(P) \gets Direct(P)$ for every procedure $P$
    \State $W \gets$ all procedures
    \While{$W \neq \emptyset$}
      \State remove some $P$ from $W$
      \State $New \gets Direct(P) \cup \bigcup_{Q \text{ called by } P} Mod(Q)$
      \If{$New \neq Mod(P)$}
        \State $Mod(P) \gets New$
        \State add every caller of $P$ to $W$
      \EndIf
    \EndWhile
  \end{pseudocode}
\end{algorithm}

The algorithm assumes direct store sets are known from a local scan and that the call graph is fixed during solving.
Represent an unknown callee by a synthetic node whose direct set contains every tracked location it may modify.
Initialize that node before iteration; omitting its edge or assigning it an empty set would make every transitive caller's summary unsound.
Recursion is handled naturally: cycles keep the worklist nonempty until the component stabilizes.
The result is the least fixed point, meaning the smallest sets satisfying all the inclusion rules, which is the most precise sound answer the rule system admits.
Clients read the summaries at call sites.
A redundant-load analysis keeps its available set across a call to $P$ except for locations in $Mod(P)$.
A dead-store analysis instead needs a may-reference summary: a read-only callee can observe a preceding store even when $Mod(P)$ is empty.
Conversely, membership in $Mod(P)$ does not prove that the call overwrites the location on every path.
An inliner uses small summaries to estimate the benefit of replacing the call with the body.

\section{Worked Example: Mod Sets Bottom-Up}
\label{sec:ch22-worked-mod}

We run Algorithm~\ref{alg:ch22-maymod} on Listing~\ref{lst:ch22-calls} by hand.
The direct store sets from scanning each body are $Direct(h) = \{g1\}$, $Direct(g) = \{g2\}$, $Direct(f) = \emptyset$, and $Direct(main) = \emptyset$.
Initialization copies these into the summaries: $Mod(h) = \{g1\}$, $Mod(g) = \{g2\}$, $Mod(f) = \emptyset$, and $Mod(main) = \emptyset$.
Processing in reverse topological order visits $h$ first.
Procedure $h$ calls nothing, so $New = \{g1\} = Mod(h)$ and nothing changes.
Next comes $g$, which calls $h$: $New = \{g2\} \cup Mod(h) = \{g1, g2\}$, which differs from $Mod(g)$, so $Mod(g)$ becomes $\{g1, g2\}$ and $g$'s caller $f$ joins the worklist.
Next comes $f$, which calls $g$: $New = \emptyset \cup Mod(g) = \{g1, g2\}$, so $Mod(f)$ becomes $\{g1, g2\}$ and $main$ joins the worklist.
Finally $main$ calls $f$: $New = \emptyset \cup Mod(f) = \{g1, g2\}$, so $Mod(main) = \{g1, g2\}$.
A final sweep changes nothing, and the fixed point is reached after one pass per procedure because the graph is acyclic.
Table~\ref{tab:ch22-mod} records the direct sets, the callee contributions, and the final summaries.

\begin{table}
  \centering
  \caption{May-modify sets for Listing~\ref{lst:ch22-calls} computed bottom-up over the call graph.}
  \label{tab:ch22-mod}
  \begin{tabular}{llll}
    \toprule
    Procedure & Direct stores & Callees add & Final summary \\
    \midrule
    $h$ & $g1$ & none & $\{g1\}$ \\
    $g$ & $g2$ & $g1$ from $h$ & $\{g1, g2\}$ \\
    $f$ & none & $g1, g2$ from $g$ & $\{g1, g2\}$ \\
    $main$ & none & $g1, g2$ from $f$ & $\{g1, g2\}$ \\
    \bottomrule
  \end{tabular}
\end{table}

The payoff is visible in \code{main}: after \code{call f()}, both globals are possibly written, so no load in \code{main} crosses the call unless its address is known distinct from both.
Had $f$ modified nothing, its empty summary would let every available load survive the call, which is exactly the precision interprocedural analysis buys.
Now introduce one complication to test understanding: if $g$ called itself recursively, the reverse-topological order would not exist.
The worklist would then process $g$ repeatedly, growing $Mod(g)$ until the self-edge contributed nothing new, and termination would still follow from monotonic growth over a finite universe.

\section{Context Sensitivity and Its Cost}
\label{sec:ch22-context}

One summary per procedure merges all call sites, and merging loses facts.
A helper called with \code{\&x} in one place and \code{\&y} in another gets a combined summary mentioning both, so neither caller recovers its precise case.
\term{Context sensitivity} keeps call sites apart.
\term{Call-site cloning} analyzes separate contexts for selected call sites, reducing caller merging without making the underlying abstraction exact.
Recursion requires a context bound or another finite representation.
\term{Summary-based sensitivity} parameterizes one summary by abstract inputs, instantiating it per call site without re-analyzing the body.
\term{Context insensitivity} accepts the merge and stays cheap.
Table~\ref{tab:ch22-sensitivity} compares the options.

\begin{table}
  \centering
  \caption{Context-sensitivity choices trade caller precision against analysis cost.}
  \label{tab:ch22-sensitivity}
  \begin{tabular}{lll}
    \toprule
    Choice & Precision & Cost driver \\
    \midrule
    Insensitive & Merged over callers & One summary per procedure \\
    Cloning & Selected contexts separated & Copies scale with contexts \\
    Summaries & Per-caller instantiation & Summary size and instantiation \\
    \bottomrule
  \end{tabular}
\end{table}

Flow sensitivity adds an orthogonal axis: a flow-sensitive analysis tracks facts per program point inside each procedure, while a flow-insensitive one merges each procedure into one fact set.
The cheapest combination, insensitive in both axes, already resolves the common cases of direct calls with global effects.
The most expensive combination is reserved for precision-critical clients like strong alias answers in security checking.
Most compilers land in the middle: flow-sensitive within procedures, context-insensitive across them, with selective cloning for tiny leaf helpers.
Separate compilation bounds every choice.
Code outside the analyzed units gets conservative summaries, and link-time analysis exists to shrink that unknown region.
% TODO(cite): primary references for context-sensitivity designs and whole-program analysis frameworks.
Soundness also requires care with callbacks, dynamic loading, and reflection-like features where the call graph itself is incomplete.
An edge the graph omits is a fact the analysis never propagates, which is unsound rather than merely imprecise.

\section{Implementation Notes}
\label{sec:ch22-implementation}

Let $P$ count procedures, $C$ call edges, and $V$ tracked locations.
A delta solver that sends each newly discovered location across each edge once takes $O(PV+CV)$ membership work with constant-time insertion.
The displayed solver recomputes whole unions instead; its cost includes repeated callee scans and $\lceil V/w\rceil$ words per bitset operation on a word size $w$.
Do not assign the delta solver's bound to a full-union implementation.
Processing in reverse topological order reaches the fixed point in one pass for acyclic graphs, and strongly connected components need iteration bounded by the height of the location lattice.
Call-graph construction for direct calls is linear in program size.
Indirect calls cost whatever the underlying pointer analysis costs, which dominates the budget in object-oriented programs with many virtual targets.
Production compilers bound context sensitivity with budgets: clone only small callees, inline the tiny ones outright, and summarize the rest.
They invalidate summaries incrementally when a procedure changes, reprocessing only its callers transitively rather than recomputing the world.
They also layer analyses by demand: compute cheap insensitive summaries for everything, then refine only the procedures on hot paths or the queries that clients actually ask.
% TODO(cite): engineering reports on scaling whole-program analysis to large codebases.
A teaching compiler should implement the acyclic worklist first, add cycle iteration with a counter that proves termination on recursive tests, and only then experiment with cloning.
Keep the unknown-code policy explicit and conservative from the start, because every later refinement inherits its soundness from that default.

\section{Exercises}
\label{sec:ch22-exercises}

\begin{enumerate}
  \item \label{ex:ch22-cycle-fixpoint}
  Extend Listing~\ref{lst:ch22-calls} so that \code{g} calls \code{f} in addition to \code{h}, creating a cycle between \code{f} and \code{g}.
  Run Algorithm~\ref{alg:ch22-maymod} with a worklist until no set changes, showing each update in order.
  Give the final summaries and explain why reverse topological order alone no longer suffices.
  \item \label{ex:ch22-merge-loss}
  Write a procedure \code{write} that stores through its pointer parameter, called once with \code{\&x} and once with \code{\&y}.
  Compute the context-insensitive may-modify summary of \code{write} and state what each caller can conclude about the other variable.
  Show the two per-caller summaries that cloning would produce, and give the load elimination in one caller that cloning enables and merging defeats.
  \item \label{ex:ch22-unknown-call}
  Add to \code{main} a call to an externally defined procedure with no visible body that returns no value.
  State the conservative may-modify summary the compiler must assume and justify it from the soundness rule for unseen code.
  Explain which loads in \code{main} survive that call under available-load analysis and which do not.
\end{enumerate}

\section{Further reading}
\label{sec:ch22-further-reading}

Safonov's \emph{Trustworthy Compilers} discusses modular reasoning and the obligations procedures owe their callers~\cite{safonov-2010-trustworthy}.
Read it for why summaries are contracts rather than mere caches.
Morgan's \emph{Building an Optimizing Compiler} connects interprocedural facts to the optimizations they enable~\cite{morgan-1998-optimizing}.
Use it to see how summaries flow into motion, numbering, and code generation.
Appel's \emph{Modern Compiler Implementation in ML} organizes these analyses inside a complete compiler~\cite{appel-1998-modern}.
Consult it when deciding where summary computation belongs in the pass pipeline.
